\documentclass[28pt, a0paper, margin=0mm,
               innermargin=4mm, blockverticalspace=0mm,
               colspace=0mm, subcolspace=4mm,
               titletotopverticalspace=4mm,
               titletoblockverticalspace=0mm]{tikzposter}

\usepackage[utf8]{inputenc}
\usepackage[T1]{fontenc}
\usepackage[czech]{babel}
\usepackage{amsmath, amssymb}
\usepackage{graphicx}
\usepackage{booktabs}

\graphicspath{{../output/}}

% ------------------------------------------------------------------
% Barevné schéma a styl bloků (ploché obdélníky, žlutá hlavička a pozadí)
% ------------------------------------------------------------------
\definecolorstyle{MSMStyle}{
    \definecolor{colorOne}{HTML}{E0B100}   % žlutá - title/block hlavičky i pozadí
    \definecolor{colorTwo}{HTML}{E0B100}
    \definecolor{colorThree}{RGB}{245,247,250}
}{
    \colorlet{backgroundcolor}{colorOne}
    \colorlet{framecolor}{colorOne}
    \colorlet{titlefgcolor}{black}
    \colorlet{titlebgcolor}{white}
    \colorlet{blocktitlebgcolor}{colorOne}
    \colorlet{blocktitlefgcolor}{black}
    \colorlet{blockbodybgcolor}{white}
    \colorlet{blockbodyfgcolor}{black}
    \colorlet{notefgcolor}{black}
    \colorlet{notefrcolor}{colorOne}
    \colorlet{notebgcolor}{colorThree}
}

\usetheme{Default}          % téma musí být PŘED vlastními styly, jinak je přepíše
\usecolorstyle{MSMStyle}
\useblockstyle[roundedcorners=0pt, titlewidthscale=1, bodywidthscale=1,
               titleinnersep=4mm, bodyinnersep=5mm]{Basic}
\usetitlestyle[roundedcorners=0pt, innersep=6mm,
               titletotopverticalspace=0mm, titletoblockverticalspace=0mm]{Filled}

\tikzposterlatexaffectionproofoff

\newcommand{\var}[1]{\texttt{\detokenize{#1}}}

\title{\parbox{0.95\linewidth}{\centering Testy o poloze: liší se chronotypy délkou spánku a únavou?\\
\large Jednorozměrné a vícerozměrné srovnání skupin Morning Lark, Intermediate a Night Owl}}
\author{Vít Novotný \& Adam Kříž}
\institute{ZČU FAV -- KMA/MSM 2026/2027}

\begin{document}
\maketitle

\begin{columns}

% ===================== SLOUPEC 1 =====================
\column{0.5}

\block{Cíl a data}{
\textbf{Otázka:} liší se tři chronotypy (\emph{Morning Lark}, \emph{Intermediate},
\emph{Night Owl}) polohou čtyř spojitých proměnných? Hladina významnosti
$\alpha = 0{,}05$ ve všech testech.

\textbf{Data:} Kaggle, \emph{Sleep Debt \& Screen Time} (syntetická data),
$n = 8\,500$ pozorování ve skupinách po $3\,880$ / $2\,165$ / $2\,455$.

\begin{minipage}[c]{0.47\linewidth}
\begin{itemize}
\item \var{bedtime_phone_minutes} -- minuty s telefonem před spaním
\item \var{sleep_latency_min} -- doba usínání (min)
\item \var{total_sleep_hours} -- délka spánku (h)
\item \var{next_day_fatigue_score} -- únava další den (0--10, strop 10)
\end{itemize}
\end{minipage}\hfill
\begin{minipage}[c]{0.51\linewidth}
\centering\small
\begin{tabular}{lrrrr}
\toprule
průměr (sm. odch.) & telefon & latence & spánek & únava \\
\midrule
Intermediate & 59,6 (37,2) & 40,8 (17,7) & 6,44 (1,17) & 3,48 (2,54) \\
Morning Lark & 59,8 (37,7) & 40,9 (17,8) & 7,38 (1,19) & 2,53 (2,18) \\
Night Owl    & 58,3 (35,9) & 40,2 (16,7) & 5,01 (1,05) & 5,40 (2,54) \\
\bottomrule
\end{tabular}
\normalsize
\end{minipage}
}

\block{Rozdělení dat}{
Telefon a latence jsou zešikmené doprava, únava má strop 10 a spánek
dolní hranici 3,2~h; normalita neplatí (Shapiro--Wilk $p<10^{-17}$ ve všech skupinách).

\centering
\includegraphics[width=\linewidth]{2_histograms.png}
}

\block{Jednovýběrové testy}{
Pro každou proměnnou a skupinu:
$H_0\!: \mu = \mu_0$ vs. $H_1\!: \mu \ne \mu_0$ (t-test) a
$H_0\!: m = m_0$ vs. $H_1\!: m \ne m_0$ (Wilcoxonův test).
Referencí jsou \textbf{průměr $\mu_0$ a medián $m_0$ všech dat}, test tedy
odpovídá na otázku, zda se skupina liší od celé populace.

\vspace{0.5em}
\begin{minipage}[c]{0.45\linewidth}
\centering\small
\begin{tabular}{llrr}
\toprule
proměnná & skupina & $p_t$ & $p_W$ \\
\midrule
telefon & Intermediate & 0,60 & $<0{,}001$ \\
telefon & Morning Lark & 0,53 & $<0{,}001$ \\
telefon & Night Owl    & 0,19 & $<0{,}001$ \\
latence & Intermediate & 0,57 & $<0{,}001$ \\
latence & Morning Lark & 0,56 & $<0{,}001$ \\
latence & Night Owl    & 0,18 & $<0{,}001$ \\
spánek  & všechny      & $<0{,}001$ & $<0{,}001$ \\
únava   & Intermediate & $<0{,}001$ & \textbf{0,93} \\
únava   & Morning Lark & $<0{,}001$ & $<0{,}001$ \\
únava   & Night Owl    & $<0{,}001$ & $<0{,}001$ \\
\bottomrule
\end{tabular}
\normalsize
\end{minipage}\hfill
\begin{minipage}[c]{0.52\linewidth}
\raggedright
Testy se rozcházejí u telefonních minut a latence: t-test průměr nezamítá,
Wilcoxon zamítá, protože průměr (tažený dlouhým pravým chvostem) padá na
$\mu_0$, ale medián ne. Opačně u únavy v~Intermediate: průměr se od $\mu_0$ liší,
medián ne.
\end{minipage}

\vspace{0.5em}
\centering
\includegraphics[width=\linewidth]{2_means_vs_reference.png}
}

\block{Dvouvýběrové testy}{
Pro každou dvojici skupin: $H_0\!: \mu_A = \mu_B$ (Welchův t-test) a
$H_0\!: F_A = F_B$ (Mann--Whitney). Levene rozhodl mezi Welchem a sdruženým
testem, obě varianty dávají shodné závěry. P-hodnoty zde \emph{nejsou}
opravené na mnohonásobné testování.

Telefon a latence se mezi chronotypy neliší ($p\ge0{,}16$). Délka spánku
a únava se liší ve \emph{všech} dvojicích ($p<0{,}001$) a oba přístupy
se shodují.
}


% ===================== SLOUPEC 2 =====================
\column{0.5}

\block{Vícevýběrové testy}{
$H_0\!: \mu_1=\mu_2=\mu_3$ (ANOVA, Welchova ANOVA) a
Kruskal--Wallis. Velikost účinku $\eta^2$ ukazuje podíl rozptylu
vysvětlený skupinou.

\vspace{0.5em}
\begin{minipage}[c]{0.5\linewidth}
\centering\small
\begin{tabular}{lrrrrr}
\toprule
proměnná & $F$ & $p_{\mathrm{ANOVA}}$ & $H$ & $p_{\mathrm{KW}}$ & $\eta^2$ \\
\midrule
telefon & 1,14 & 0,32 & 0,96 & 0,62 & 0,0003 \\
latence & 1,17 & 0,31 & 0,53 & 0,77 & 0,0003 \\
spánek  & 2551 & $<0{,}001$ & 3244 & $<0{,}001$ & 0,38 \\
únava   & 845  & $<0{,}001$ & 1619 & $<0{,}001$ & 0,17 \\
\bottomrule
\end{tabular}
\normalsize
\end{minipage}\hfill
\begin{minipage}[c]{0.47\linewidth}
\raggedright
Post-hoc (Tukey HSD, Mann--Whitney s Holmovou korekcí) dává pro všech
12 dvojic \emph{stejné rozhodnutí} jako nekorigované testy z předchozího
bloku. Při $n=8\,500$ jsou rozdíly tak velké, že korekce nic nezmění.
\end{minipage}

\vspace{0.5em}
\centering
\includegraphics[width=\linewidth]{3_boxplots.png}
}

\block{Diagnostika předpokladů}{
Rozptyly se podle Leveneova a Bartlettova testu liší u spánku a únavy
($p<0{,}001$), rezidua nejsou normální (Shapiro na vzorku 5\,000).
Welchova ANOVA a Kruskal--Wallis proto dávají spolehlivější závěry než
klasická ANOVA, výsledky ale vycházejí stejně.
}

\block{Vícerozměrný přístup}{
Vektor $(\text{telefon},\text{latence},\text{spánek},\text{únava})^\top$.
Jednovýběrový $T^2$ ($H_0\!: \boldsymbol\mu=\boldsymbol\mu_0$) zamítá
pro všechny skupiny ($T^2 = 259$ / $4\,411$ / $7\,246$), dvouvýběrový $T^2$
zamítá pro všechny dvojice i po Holmově korekci
($T^2 = 1\,999$ / $5\,217$ / $10\,820$). \textbf{MANOVA}: Wilksovo
$\Lambda = 0{,}429$, Pillaiova stopa $0{,}578$ ($F(8;16\,990)=863$, $p<0{,}001$).

Boxův $M$ test zamítá shodu kovariančních matic ($\chi^2(20)=784$,
$p<0{,}001$), proto bereme MANOVA a $T^2$ jako orientační; Pillaiova stopa
je vůči porušení nejodolnější.

\vspace{0.5em}
\centering
\includegraphics[width=0.8\linewidth]{5_mean_profiles.png}

\includegraphics[width=0.8\linewidth]{5_t2_difference_ci.png}

\begin{minipage}[c]{0.4\linewidth}
\centering
\includegraphics[width=\linewidth]{5_centroid_distances.png}
\end{minipage}\hfill
\begin{minipage}[c]{0.57\linewidth}
\raggedright
Mahalanobisova vzdálenost průměrových vektorů: nejdál od sebe jsou
Morning Lark a Night Owl ($D^2 = 9{,}41$), Intermediate leží mezi nimi.
Simultánní $T^2$ intervaly (modře) jsou širší než jednotlivé t-intervaly,
závěry se ale nemění.
\end{minipage}
}

\block{Závěr}{
\begin{itemize}
\item Chronotypy se liší \textbf{délkou spánku} (Night Owl $\approx 5{,}0$ h,
      Morning Lark $\approx 7{,}4$ h) a \textbf{únavou}, ne telefonem před
      spaním ani latencí usínání.
\item Parametrické a neparametrické přístupy se shodují ve všech
      dvojicích; rozdílné jsou jen jednovýběrové testy u zešikmených
      proměnných (průměr vs. medián).
\item Při $n=8\,500$ je test citlivý i na malé rozdíly, proto uvádíme
      $\eta^2$: spánek $0{,}38$, únava $0{,}17$, ostatní $<0{,}001$.
\item Normalita ani shoda kovariančních matic neplatí; závěry jsou
      robustní díky velkému $n$ a velikosti rozdílů.
\item Vícerozměrný pohled potvrzuje závěry po proměnných a ukazuje, že
      skupiny se liší hlavně v kombinaci spánku a únavy.
\end{itemize}
}

\block{Technologie a použití AI}{
Python 3.13, NumPy, pandas, SciPy, statsmodels, Matplotlib. \texttt{main.py} provede všechny testy a uloží tabulky a grafy do \texttt{output/}; $T^2$ a Boxův $M$ test jsou implementovány ručně.

\textbf{Použití AI:} V rámci této práce byla využita umělá inteligence (Claude a Gemini) k odstraňování chyb a optimalizaci kódu v jazyce Python.
}

\end{columns}

\end{document}
